% Source: https://artofproblemsolving.com/community/c3h1656531_2018_ms_mte_sprint_all
% Stable identity: aops-topic-1656531-q07
Use coordinates $A=(0,0)$, $B=(1,0)$, $C=(1,1)$, $D=(0,1)$. Then $M=\left(\dfrac{1}{2},0\right)$, $T=\left(1,\dfrac{2}{3}\right)$ (as $TB=\dfrac{2}{3}$) and $H=\left(\dfrac{1}{4},1\right)$ (as $HD=\dfrac{1}{4}$).

\begin{asy}
usepackage("amsmath");
size(5cm);
defaultpen(fontsize(10pt));
pair A=(0,0),B=(1,0),C=(1,1),D=(0,1);
pair M=(0.5,0),T=(1,2/3),H=(0.25,1);
fill(A--M--T--H--cycle,palegray);
draw(A--B--C--D--cycle,linewidth(0.9));draw(A--M--T--H--cycle,blue+linewidth(0.9));
dot(M);dot(T);dot(H);
label("$A$",A,SW);label("$B$",B,SE);label("$C$",C,NE);label("$D$",D,NW);
label("$M$",M,S);label("$T$",T,E);label("$H$",H,N);
label("\mbox{}",(min(currentpicture,user=true)+max(currentpicture,user=true))/2,black); // keeps the text after the figure black when LaTeX typesets the labels inline
\end{asy}

The convex quadrilateral is $AMTH$. Subtract the three corner pieces from the square:
\[
[MBT]=\dfrac{1}{2}\cdot\dfrac{1}{2}\cdot\dfrac{2}{3}=\dfrac{1}{6},\qquad [TCH]=\dfrac{1}{2}\cdot\dfrac{1}{3}\cdot\dfrac{3}{4}=\dfrac{1}{8},\qquad [HDA]=\dfrac{1}{2}\cdot\dfrac{1}{4}\cdot1=\dfrac{1}{8} .
\]
So $[AMTH]=1-\dfrac{1}{6}-\dfrac{1}{8}-\dfrac{1}{8}=\boxed{\textbf{(c)}\ \dfrac{7}{12}}$.
